%%
%% This is file `sigconf.tex',
%% generated with the docstrip utility.
%%
%% The original source files were:
%%
%% samples.dtx  (with options: `all,proceedings,sigconf')
%% 
%% IMPORTANT NOTICE:
%% 
%% For the copyright see the source file.
%% 
%% Any modified versions of this file must be renamed
%% with new filenames distinct from sigconf.tex.
%% 
%% For distribution of the original source see the terms
%% for copying and modification in the file samples.dtx.
%% 
%% This generated file may be distributed as long as the
%% original source files, as listed above, are part of the
%% same distribution. (The sources need not necessarily be
%% in the same archive or directory.)
%%
%%
%% Commands for TeXCount
%TC:macro \cite [option:text,text]
%TC:macro \citep [option:text,text]
%TC:macro \citet [option:text,text]
%TC:envir table 0 1
%TC:envir table* 0 1
%TC:envir tabular [ignore] word
%TC:envir displaymath 0 word
%TC:envir math 0 word
%TC:envir comment 0 0
%%
%% The first command in your LaTeX source must be the \documentclass
%% command.
%%
%% For submission and review of your manuscript please change the
%% command to \documentclass[manuscript, screen, review]{acmart}.
%%
%% When submitting camera ready or to TAPS, please change the command
%% to \documentclass[sigconf]{acmart} or whichever template is required
%% for your publication.
%%
%%
% \documentclass[manuscript, screen, review]{acmart}  % For submit to review
% \documentclass[sigconf]{acmart}
% \documentclass[sigconf] {acmart}
\documentclass[sigconf, nonacm] {acmart}

\usepackage{xcolor}

\usepackage{tikz}



%%
%% \BibTeX command to typeset BibTeX logo in the docs
\AtBeginDocument{%
  \providecommand\BibTeX{{%
    Bib\TeX}}}

%% Rights management information.  This information is sent to you
%% when you complete the rights form.  These commands have SAMPLE
%% values in them; it is your responsibility as an author to replace
%% the commands and values with those provided to you when you
%% complete the rights form.

% \setcopyright{acmlicensed}
% \copyrightyear{20XX}
% \acmYear{20XX}
% \acmDOI{XXXXXXX.XXXXXXX}


%% These commands are for a PROCEEDINGS abstract or paper.

% \acmConference[Conference acronym 'XX]{Make sure to enter the correct
%   conference title from your rights confirmation email}{June 03--05,
%   2018}{Woodstock, NY}


%%
%%  Uncomment \acmBooktitle if the title of the proceedings is different
%%  from ``Proceedings of ...''!
%%
%%\acmBooktitle{Woodstock '18: ACM Symposium on Neural Gaze Detection,
%%  June 03--05, 2018, Woodstock, NY}

% \acmISBN{978-1-4503-XXXX-X/2018/06}


%%
%% Submission ID.
%% Use this when submitting an article to a sponsored event. You'll
%% receive a unique submission ID from the organizers
%% of the event, and this ID should be used as the parameter to this command.
%%\acmSubmissionID{123-A56-BU3}

%%
%% For managing citations, it is recommended to use bibliography
%% files in BibTeX format.
%%
%% You can then either use BibTeX with the ACM-Reference-Format style,
%% or BibLaTeX with the acmnumeric or acmauthoryear sytles, that include
%% support for advanced citation of software artefact from the
%% biblatex-software package, also separately available on CTAN.
%%
%% Look at the sample-*-biblatex.tex files for templates showcasing
%% the biblatex styles.
%%
%%
%% The majority of ACM publications use numbered citations and
%% references.  The command \citestyle{authoryear} switches to the
%% "author year" style.
%%
%% If you are preparing content for an event
%% sponsored by ACM SIGGRAPH, you must use the "author year" style of
%% citations and references.
%% Uncommenting
%% the next command will enable that style.
%%\citestyle{acmauthoryear}


%%
%% end of the preamble, start of the body of the document source.
\begin{document}

%%
%% The "title" command has an optional parameter,
%% allowing the author to define a "short title" to be used in page headers.
\title{Reliability First: Lessons from Industrial HRI Deployment in High-Stakes Environments}

  % When Robot Deployments Fail: Evidence, Mechanisms, and Lessons for Industry Practice
% Toward a Taxonomy of HRI Deployment Failures: Evidence from Real-World Cases and Lessons for Industry Practice


%%
%% The "author" command and its associated commands are used to define
%% the authors and their affiliations.
%% Of note is the shared affiliation of the first two authors, and the
%% "authornote" and "authornotemark" commands
%% used to denote shared contribution to the research.
% \author{Ben Trovato}
% \authornote{Both authors contributed equally to this research.}
% \email{trovato@corporation.com}
% \orcid{1234-5678-9012}
% \author{G.K.M. Tobin}
% \correspondingauthor
% \authornotemark[1]
% \email{webmaster@marysville-ohio.com}
% \affiliation{%
%   \institution{Institute for Clarity in Documentation}
%   \city{Dublin}
%   \state{Ohio}
%   \country{USA}
% }


\author{Arvind}
\orcid{0000-0002-0456-7122}
\affiliation{%
  \institution{Academy of Scientific \\ and Innovative Research (AcSIR),}
  % \city{Ghaziabad}
  \country{India}
}
\affiliation{%
  \institution{CSIR-Central Electronics Engineering \\ Research Institute (CEERI), Pilani}
  % \city{Pilani}
  \country{India}
}
\email{arvind.ceeri24a@acsir.res.in}

\author{Sumeet Saurav}
\orcid{0000-0002-4375-4107}
\affiliation{%
  \institution{Academy of Scientific \\ and Innovative Research (AcSIR),}
  % \city{Ghaziabad}
  \country{India}
}
\affiliation{%
  \institution{CSIR-Central Electronics Engineering \\ Research Institute (CEERI), Pilani}
  % \city{Pilani}
  \country{India}
}
\email{sumeet@ceeri.res.in}

% \author{Sanjay Singh}
% \orcid{0000-0002-2249-799X}
% \affiliation{%
%   \institution{Academy of Scientific \\ and Innovative Research (AcSIR),}
%   % \city{Ghaziabad}
%   \country{India}
% }
% \affiliation{%
%   \institution{CSIR-Central Electronics Engineering \\ Research Institute (CEERI), Pilani}
%   % \city{Pilani}
%   \country{India}
% }
% \email{sanjay@ceeri.res.in}




% \author{Arvind}
% \affiliation{%
%   \institution{The Th{\o}rv{\"a}ld Group}
%   \city{Hekla}
%   \country{Iceland}}
% \email{larst@affiliation.org}

% \author{Valerie B\'eranger}
% \affiliation{%
%   \institution{Inria Paris-Rocquencourt}
%   \city{Rocquencourt}
%   \country{France}
% }

% \author{Aparna Patel}
% \affiliation{%
%  \institution{Rajiv Gandhi University}
%  \city{Doimukh}
%  \state{Arunachal Pradesh}
%  \country{India}}

% \author{Huifen Chan}
% \affiliation{%
%   \institution{Tsinghua University}
%   \city{Haidian Qu}
%   \state{Beijing Shi}
%   \country{China}}

% \author{Charles Palmer}
% \affiliation{%
%   \institution{Palmer Research Laboratories}
%   \city{San Antonio}
%   \state{Texas}
%   \country{USA}}
% \email{cpalmer@prl.com}

% \author{John Smith}
% \affiliation{%
%   \institution{The Th{\o}rv{\"a}ld Group}
%   \city{Hekla}
%   \country{Iceland}}
% \email{jsmith@affiliation.org}

% \author{Julius P. Kumquat}
% \correspondingauthor
% \affiliation{%
%   \institution{The Kumquat Consortium}
%   \city{New York}
%   \country{USA}}
% \email{jpkumquat@consortium.net}

%%
%% By default, the full list of authors will be used in the page
%% headers. Often, this list is too long, and will overlap
%% other information printed in the page headers. This command allows
%% the author to define a more concise list
%% of authors' names for this purpose.
% \renewcommand{\shortauthors}{Trovato et al.}

\renewcommand{\shortauthors}{Arvind et al.}

%%
%% The abstract is a short summary of the work to be presented in the
%% article.
\begin{abstract}
Industry deployments of human-robot interaction (HRI) systems in high-stakes areas like inspection, maintenance, and precision manufacturing shows a continuing conflict between the introduction of new interaction methods and the need for operational reliability. This white paper draws on field evidence from recent industrial deployments of HRI systems in order to highlight patterns that affect adoption, user acceptance, and long-term sustainability. We claim that in situations which are conservative and cautious from the point of view of operations, constraints relating to reliability consistently prevent the introduction of new interaction features, that multi-operator workflows stay the rule rather than the exception, and that there is a gap between aspirational usability and reality which damages the trust between system designers and the end users. Based on survey data, field studies, and deployment reports, we have identified four lessons that can be applied by HRI practitioners and researchers: (1) design for the operator team, not for the individual; (2) regard trust as a property that varies depending on the stakeholders involved and on the task in question; (3) bridge the gap between intended and actual users; and (4) realise that complying with standards does not mean that a system is easy to use. These lessons are not intended to be treated as rigid rules but rather as diagnostic patterns to assist both industrial practitioners and academic researchers in anticipating the socio-technical difficulties that occur in real-world HRI situations.
Supporting material available at: \href{https://github.com/arvindsihag/Decobots}{https://github.com/arvindsihag/Decobots}.


\end{abstract}

%%
%% The code below is generated by the tool at http://dl.acm.org/ccs.cfm.
%% Please copy and paste the code instead of the example below.
%%
% \begin{CCSXML}
% <ccs2012>
%  <concept>
%   <concept_id>00000000.0000000.0000000</concept_id>
%   <concept_desc>Do Not Use This Code, Generate the Correct Terms for Your Paper</concept_desc>
%   <concept_significance>500</concept_significance>
%  </concept>
%  <concept>
%   <concept_id>00000000.00000000.00000000</concept_id>
%   <concept_desc>Do Not Use This Code, Generate the Correct Terms for Your Paper</concept_desc>
%   <concept_significance>300</concept_significance>
%  </concept>
%  <concept>
%   <concept_id>00000000.00000000.00000000</concept_id>
%   <concept_desc>Do Not Use This Code, Generate the Correct Terms for Your Paper</concept_desc>
%   <concept_significance>100</concept_significance>
%  </concept>
%  <concept>
%   <concept_id>00000000.00000000.00000000</concept_id>
%   <concept_desc>Do Not Use This Code, Generate the Correct Terms for Your Paper</concept_desc>
%   <concept_significance>100</concept_significance>
%  </concept>
% </ccs2012>
% \end{CCSXML}

% \ccsdesc[500]{Do Not Use This Code~Generate the Correct Terms for Your Paper}
% \ccsdesc[300]{Do Not Use This Code~Generate the Correct Terms for Your Paper}
% \ccsdesc{Do Not Use This Code~Generate the Correct Terms for Your Paper}
% \ccsdesc[100]{Do Not Use This Code~Generate the Correct Terms for Your Paper}

%%
%% Keywords. The author(s) should pick words that accurately describe
%% the work being presented. Separate the keywords with commas.
\keywords{human-robot interaction, industry deployment, reliability, trust, multi-operator workflows, field studies}

\maketitle


\section{Introduction}
Over the last twenty years, research into human-robot interaction has advanced rapidly, Hancock et al.'s meta-analysis found that trust in robots is influenced by performance, process, and purpose factors~\cite{hancock2021evolving}. This is consistent with the results of more general trust research. Trust in automation can be understood in terms of performance, process, and purpose factors~\cite{lee2004trust}, collaboration~\cite{nikolaidis2017human}, and communication~\cite{arkin2018realtime}. 


However, there is still a considerable gap between the demonstrations carried out in laboratories and actual, long-term industrial applications. According to recent surveys of people making technology-related decisions, although 71\% of organisations currently use robotics or intend to do so, trust in this area remains weak and depends on having demonstrated safety and reliability~\cite{QNX2025,QNX2025de}. Almost one-third of the respondents from Germany said that they felt uncomfortable working with robots, and 41\% mentioned the high initial costs as a reason why they were not doing so~\cite{QNX2025,QNX2026InsideTheRobot,QNX2026ActiveSafety,QNX2026FunctionalSafety}.

This white paper focuses on one aspect of the gap referred to, namely the use of HRI systems in high-stakes operational environments in which a failure of the system has material, safety, or regulatory implications. Based on recent field evidence obtained from applications in manufacturing, inspection, and collaborative robotics, we have identified four patterns that repeatedly occur and which influence whether or not HRI systems succeed or fail in real-world situations. The contribution we are making is not the presentation of any new empirical data but rather a synthesis of field lessons that have already been documented, presented in a way that is useful to both practitioners in the industry and to academic researchers.

The paper is structured as follows: Section II sets out the reliability imperative which determines the design decisions in high-stakes areas. Section III looks at the challenge associated with the multi-operator workflow. Section IV investigates the trust dynamics that exist among the different stakeholder roles. Section V summarizes four lessons that can be applied in other situations. Section VI concludes.

\section{The Reliability Imperative}
In situations where the stakes are high, the most important design consideration is not the user's experience or the elegance of the interaction; it is rather operational reliability. This is evident in a number of ways.

In the first place, the interaction methods remain intentionally conservative. Industrial teleoperation interfaces usually make use of physical joysticks and button panels rather than gesture, voice, or touchscreen control, a practice being justified by the experience gained by operators and by the requirement for fail-safe operation in degraded conditions. As a field study concerning the control of bimanual robotic arms noted, the main performance criterion is not intuitive interaction but rather "movement speed, time on task, or accuracy" within the context of production constraints~\cite{Dreger02062024}.

Second, autonomy is very carefully restricted. In pilot projects involving vision systems for food production, extensive recalibration was necessary because of interference from ambient light before the systems could operate reliably~\cite{hartu2025field}. 

The same kind of situation occurs in collaborative manufacturing. Gervasi et al.\ showed that robot movement speed and control of task execution time can significantly influence user experience, affective state, and physiological stress during collaborative assembly tasks~\cite{Gervasi2022}. 

Related experimental evidence also shows that assembly complexity and cobot assistance influence physiological indicators
of cognitive effort, motivating multidimensional evaluation of worker well-being~\cite{CAPPONI2024102789}.

The outcome is a situation in which advanced autonomous capabilities exist alongside conservative human-in-the-loop backup options.

Third, business models place greater emphasis on operational expertise than on interaction innovation. In the case of inspection robotics, the value offered is based on accumulated knowledge of the domain and on the reliability of the service rather than on innovative interaction design. As a result, a design culture develops in which "it doesn't matter how you carry out the work, just get on and do it reliably" becomes an unspoken yet widespread engineering principle.

The point for researchers in the field of human-robot interaction is clear: in high-stakes areas, any new form of interaction has to prove that it is as reliable as the existing methods before it can replace them. This is not an opposition to innovation but rather a reasonable reaction to the unequal costs of failure.




\section{The Multi-Operator Challenge}
Another frequent pattern is the continued use of multi-operator work processes. The usual promise of robotics—'reducing the number of personnel'—does not generally come about in reality since safety regulations, the need to separate expertise, and the complexity of the tasks all call for human teams rather than a single operator.

Documents relating to inspections and maintenance using robotics show that there are workflows in which different roles (for example, robot operator as compared to surveyor or inspector) have conflicting requirements when it comes to the interface. The robot operator needs real-time telemetry, camera feeds, and control rights; the surveyor on the other hand needs high-resolution inspection data, annotation tools, and analysis interfaces. When efforts are made to combine these roles into one interface, the results are generally unsatisfactory for both parties.

The same kind of situation occurs in collaborative manufacturing. Gervasi et al.\ showed that robot movement speed and control of task execution time can significantly influence user experience, affective state, and physiological stress during collaborative assembly tasks~\cite{Gervasi2022}. This reinforces the need to evaluate HRC as a multidimensional system involving human, robot, interaction, and task-related factors rather than considering robot performance alone~\cite{Gervasi2020}.

The idea is that HRI systems have to enable coordination between different
roles, not just allow interaction with one user. Recent in-the-wild studies
of group--robot interaction similarly show that interaction with robots can
become a multi-party phenomenon rather than a simple human--robot dyad~\cite{zibetti2026field}. Multi-operator interfaces should therefore
be able to meet differing information needs, deal with unequal levels of
authority, and include clear handoff procedures. This constitutes a
socio-technical design issue, not just a UI issue.





\section{Trust Across Stakeholder Roles}
There is no single, uniform level of trust in industrial human–robot interaction. Evidence from the field shows that there are systematic differences depending on the role of the stakeholder.

The people who operate robots and therefore interact most directly with the system are generally enthusiastic supporters of the autonomy features since they get to see for themselves what the system can do and as a result build up in their minds a clear understanding of how reliable it is. On the other hand, downstream users of the data—such as surveyors, inspectors, and quality engineers—are often unwilling to trust the outputs of the autonomous system because they have professional responsibility for decisions based on those outputs even though they do not have direct experience of the system's failure modes.

This is consistent with the results of more general trust research. Hancock et al.'s meta-analysis found that trust in robots is influenced by performance, process, and purpose factors, different stakeholders assigning different weights to these aspects~\cite{hancock2021evolving}. Likewise, the data from the QNX survey indicate that trust depends on the task: 77\% of the respondents were comfortable with robots carrying out assembly-line work, but only 63\% were comfortable with them carrying out maintenance and repairs and 51\% with medical procedures~\cite{qnx2025study}.

Another example is what we call the aspirational usability gap. Although marketing and procurement statements usually assert that "anyone can use" the system, in practice it is often necessary to have "a full-stack software engineer" in order to carry out effective troubleshooting and adaptation. Trust disappears whenever real users do not meet the standards set out in the aspirational user model—not since the system is unreliable, but because the difference between the promised and the actual usability leads to unfulfilled expectations.

The relevant frameworks are provided by the HRI literature. Chen et al. have created trust-aware decision-making models which take into account human trust dynamics~\cite{chen2018trustawaredecisionmakinghumanrobot}. Our field observations indicate that these models should be extended to cover multi-stakeholder trust, since different roles have different trust calibration requirements and different consequences when trust is misplaced.

\section{Four Transferable Lessons}
We are able to identify four lessons from the patterns described above which are intended to apply in a wide range of high-stakes industrial areas.

\subsection{Lesson 1: Design with the team, not for the individual}
In industrial human-machine interaction the typical unit of interaction is a team consisting of members who have different roles, levels of authority and specific information needs. Interaction models based on a single user (which are common in HRI research) can be misleading in such situations. Before deciding on the interface architecture, the design process should clearly identify the information requirements that vary according to role, the boundaries of authority and the coordination procedures.

\subsection{Lesson 2: Treat Trust as Task-Contingent and Role-Dependent}
The degree to which people trust industrial robots is not a simple scalar attribute of the system but rather a multidimensional construct that depends on the role of the individual, the task at hand, and the circumstances. Measures which aim at increasing operators' trust—such as providing transparent telemetry or offering explainable autonomy—might not lead to an increase in, and could even result in a decrease in, downstream consumer trust if they highlight uncertainty that the consumers cannot act upon. Therefore, trust-building approaches have to take into account the role of the individual.

\subsection{Lesson 3: Close the Aspirational--Actual User Gap}
The difference between the intended user ("anyone can operate") and the real user (usually a skilled technician) is a major cause of a loss of trust and difficulties in putting the system into use. In order to achieve successful HRI design it is essential to have an honest understanding of who the actual users will be—that is, who will be responsible for operating, maintaining and relying on the system? When the profile of the actual user differs from that of the intended one, training schemes and support systems have to be redesigned, not just improved.

\subsection{Lesson 4: Standards Compliance $\neq$ Usability}
The safety guidelines given in ISO 10218 and ISO/TS 15066 are important, but merely meeting them does not result in usable or acceptable human-robot interaction. These standards limit the speed and force of collaborative robots when they are near humans, which can cause the robots to seem "incompetent and hazardous" in real-world use~\cite{robotics2025}. Therefore, designers of human-robot interaction should consider the standards as a minimum requirement rather than as an upper limit and should actively focus on ensuring usability while working within the safety constraints.


\section{Conclusion}
The white paper has drawn together field evidence from industrial deployments involving human–robot interaction in order to identify four patterns that recur and influence adoption in high-stakes situations: the reliability imperative, the multi-operator challenge, role-dependent trust dynamics, and the aspirational–actual user gap. However, these patterns are not presented as absolute laws but rather as heuristic tools for use by practitioners and researchers as they move from laboratory prototypes to actual deployed systems.

The main point is that industrial HRI is not just "HRI in a factory". The constraints, the way stakeholders are organised, and the cost of failures in operational environments lead to a particular design space which calls for both technical expertise and socio-technical sensitivity. We aim to help promote a more realistic and productive discussion between industry practitioners and the academic HRI community.



\section*{Acknowledgments}
The authors thank the HRI 2027 Industry White Paper track organizers and reviewers.


\bibliographystyle{ACM-Reference-Format}
\bibliography{ref}

\end{document}